\section{Discussion}

The consequences of low-quality organoleptic assessments have been described
most thoroughly for wine. In studies of wine judges, only about 10\% of them
were able to replicate their own rating of the same blind wine sample within
a single category \cite{hodgson2008judge}. Moreover, winning a high award at
one competition did not predict the result at another: the two results were
statistically independent \cite{hodgson2009concordance}. Organoleptic
assessment can be calibrated through the selection and training of assessors,
monitoring of panel performance indicators (repeatability, reproducibility,
agreement), and the use of reference standards and externally anchored scales
\cite{sipos2021panel,bartoshuk2004valid}. For honey, such a methodology was
developed within the International Honey Commission \cite{Piana2004}, and
since 1999 Italy has maintained a national register of experts in the
organoleptic analysis of honey, entry into which requires completed training
and a passed examination \cite{AlboMiele1999}. In France, the methodology for
the organoleptic assessment of unifloral honeys was developed by the team of
M.~Gonnet (INRA, now INRAE), whose work is the origin of modern organoleptic
analysis of honey \cite{Piana2004}. In Germany, the principles of the
organoleptic assessment of honey were drawn up by the German Agricultural
Society (DLG): its stated aims include the characterization of unifloral
honeys together with pollen analysis and the detection of the effects of
overheating and prolonged storage, and the assessment scheme lists, among
others, the odor and taste defects ``not typical for honey'', ``foreign
odor'' and ``fermenting'' \cite{DLG2015}. In both countries, organoleptic
assessment complements laboratory testing rather than replacing it. In routine
honey assessment, however, these indicators are usually not published in the
reports, and laboratories do not share a common reference material. The
comparability of organoleptic verdicts between laboratories thus remains an
assumption rather than a demonstrated property of the measurement.

The same pattern is evident between laboratories. Agreement among the four
laboratories was fair on the binary scale and slight on the three-level scale
\cite{LandisKoch1977}, and full agreement concerned only samples to which no
laboratory raised objections. Samples P1 and P2 were flagged by only one of
the four laboratories, while the others described them as typical, slightly
typical or of ambiguous characteristics. Lab4, which used a different method,
issued the verdict ``ambiguous characteristics'' for all seven samples,
including P6 and P7, which were flagged by the other three laboratories, and
reported an odor characteristic of signs of fermentation in three samples,
including P5, to which none of the other laboratories raised objections. The
verdict for the same batch therefore depends not only on random variability of
perception but also on the method applied and on the laboratory.

These discrepancies do not prove an error by any particular laboratory or a
defect of any particular batch. They do show, however, that in commercial
quality control practice the typicality of odor and taste is not measured
reproducibly between laboratories, in line with the predictions of the
psychometric literature
\cite{bartoshuk2003labeled,hodgson2009concordance,sipos2021panel}. The label
``not typical'' results from a comparison with an internal, implicit reference
of the assessor, and this reference differs between persons and laboratories.

\subsection{Describing an Attribute Versus Its Decision Weight}

The panel registered an attribute present in some of the laboratory
descriptions: the sweetness of P6 and P7 was rated as ``too strong'' more
often than that of other samples, and these samples received the lowest
typicality. This relationship was not unequivocal, however. Consumers
attributed the same sweetness intensity to sample~P5, to which the laboratories
using a typicality qualifier raised no objections. Sample~P2, in turn,
described by one of the laboratories as ``very sweet'', was rarely rated by
consumers as too sweet.

These results suggest that an organoleptic description may capture an
attribute of the sample accurately and still fail to distinguish typical from
not typical honey. A possible explanation is that typicality is determined not
by the intensity of a single attribute but by its consistency with the profile
of a given variety; this, however, requires further research. It is therefore
useful to distinguish an \emph{accurate description} from a \emph{decisive
criterion}: an attribute may be useful for characterizing a product and at the
same time unsuitable as a stand-alone basis for an administrative decision
\cite{Gajda2026}.

\subsection{Implications for Quality Control}

From the point of view of control, what matters is not whether taste
assessment carries any information at all, but whether it allows a verdict to
be reproduced. A single panelist rating reproduced the majority verdict of the
laboratories with $\mathrm{AUC} = 0.57$ (95\% CI: 0.54--0.60), and a rule
based on taste typicality would be wrong in nearly every second rating. The
agreement of a consumer with a laboratory ($\kappa = 0.10$) did not differ
from the agreement of two consumers ($\kappa = 0.08$). The laboratories
themselves fully agreed on three of the seven samples on the binary scale and
on one on the three-level scale. A criterion whose result for borderline
samples depends on the choice of assessor or laboratory does not satisfy the
condition of reproducibility and, when used on its own, the condition of
proportionality.

This difference is visible in the reports themselves. Lab2 and Lab4 described
a decision rule that takes measurement uncertainty into account
\cite[clause~7.8.6]{ISO17025} with respect to the limit values of measurable
parameters, and Lab4 explicitly named the physicochemical parameters and
pollen analysis. For the organoleptic assessment, no report stated an
uncertainty or a decision rule. For instrumental methods, precision is
determined in interlaboratory studies and is used to set critical differences
applied when assessing compliance with honey purity criteria
\cite{Ulberth2024}. For the verdict ``atypical taste'', no analogous
quantities are known.

This does not mean that organoleptic assessment is useless. It is a valuable
tool for studying preferences and for describing and improving products, and
also a signal for further testing, in the role foreseen for the organoleptic
test in the coordinated control plan of the European Commission
\cite{KE2015PlanKontroli}, when the result is extreme and consistent between
independent assessors, as for P6 and P7. A similar relationship has been
described for trained experts: in the International Honey Commission trials,
the results of assessing conformity with the profile of a unifloral honey were
valid for clearly good or clearly defective samples, whereas for samples with
intermediate characteristics the spread of ratings was largest, and the binary
``yes/no'' classification gave no valid result in 23\% of cases even after
three panels had been combined \cite{Piana2004}. The results indicate that for
borderline samples an organoleptic verdict should not be confirmed by a further
organoleptic assessment conducted without a reference. A second laboratory
using its own descriptors provides another description, not a decision: one
may describe the taste as ``very sweet'', another the odor as ``candy-like'',
and without a defined reference set of typicality descriptors neither
description establishes atypicality. The verdict ``atypical taste'' should
therefore be verified with reliable results, namely physicochemical and
microscopic determinations and, where adulteration is suspected, analytical
methods of known precision \cite{Ulberth2024}, before it becomes the basis for
a decision to withdraw a product. The role of organoleptic assessment in
commercial quality control should be limited to detecting the atypical
characteristics defined in the legislation, such as a foreign taste or odor
and signs of fermentation \cite{Dyrektywa2001110,RozporzadzenieMRiRW2009},
and should not extend to judging typicality, for which no reference exists.

Such a limitation corresponds to the structure of EU law. Directive
2001/110/EC formulates the organoleptic requirements as prohibitions: honey,
with the exception of baker's honey, must not have any foreign taste or odor
and must not show signs of fermentation \cite{Dyrektywa2001110}. An assessment
of appearance, odor, taste, consistency, crystallization and visible signs of
fermentation thus makes it possible to indicate a possible non-compliance and
to justify referring the sample for further laboratory testing. It was in this
preliminary and complementary role, and not as a stand-alone analytical
method, that the organoleptic test was used in the coordinated control plan of
the European Commission \cite{KE2015PlanKontroli}.

\subsection{Limitations of the Study}
\label{sub:limitations}

The main limitation is the small number of batches. Although the panel
provided 966~ratings, the sample-level comparisons (panel--laboratory
correlation, agreement of verdicts) are based on seven units. Second, the
panel consisted of consumers after a short briefing rather than of selected
assessors; this reflects the market perspective but increases the spread of
ratings. The panel question on taste typicality is also not identical to the
laboratory procedures, which assess odor and taste jointly according to their
own references. Third, the laboratory assessment was carried out as commercial
orders rather than as a formal proficiency test: the samples were not coded,
the laboratories had no common reference material, and the coding of their
descriptions on the T/L/N scale, including the assignment of code L to the
verdict ``ambiguous characteristics'', is a simplification made by the
authors. Fourth, the panel (May--July 2026) and the laboratory assessment
(August 2026) were separated by one to three months, during which
storage-dependent characteristics (crystallization, fermentation) may have
changed. Finally, only seven batches were examined, so the results should not
be generalized to all honeys.

\section{Conclusions}

The assessment of honey taste and odor typicality conducted as in this study
(by a blind consumer panel or by a single laboratory, without a common
reference standard and without a reported measure of assessor agreement) is
too variable to serve as a decisive parameter in the assessment of the
commercial quality of a sample.

Three independent findings support this conclusion. First, the discriminative
ability of the assessment is unacceptably low: a single panelist rating
reproduced the laboratory verdict with $\mathrm{AUC} = 0.57$, i.e., only
slightly better than a random binary choice ($\mathrm{AUC} = 0.5$) and clearly
below the value of 0.7 regarded as the minimum for acceptable discrimination
\cite{Hosmer2013}. In practice, the rule ``atypical taste, sample rejected''
would reject 39.7\% of the ratings of samples to which the majority of
laboratories raised no objections, and would pass 46.4\% of the ratings of
samples flagged by the majority, so it would be wrong in nearly every second
decision. Second, apart from one clearly typical sample, the panel did not
discriminate between samples and, as a whole, flagged none of them: the
typicality of the samples flagged by most laboratories (0.438 and 0.489) was
indistinguishable from 0.5. Third, the assessors agreed neither among
themselves nor with the laboratories: agreement among the four laboratories was
fair on the binary scale ($\kappa = 0.30$) and slight on the three-level scale
($\kappa = 0.13$), and a single consumer agreed with a laboratory to a similar
degree as with another consumer ($\kappa = 0.10$ and $\kappa = 0.08$).

A decision to accept a batch as compliant or to withdraw it, based solely on
such a result, would depend more on who assesses the sample and where than on
the properties of the sample itself. A subjective test should therefore not
serve as the basis for a binary ``good/bad'' judgment. Organoleptic assessment
retains its value as a descriptive tool and as a signal for further testing,
especially when the result is extreme and consistent between independent
assessors. It should not, however, decide on its own whether honey complies
with the commercial quality requirements. In the light of these results, the
distinction between identifying actual quality non-compliance and the
subjective evaluation of the organoleptic properties of honey is of central
importance. The absence of clearly defined and objective typicality criteria
limits the possibility of using the individual impressions of an assessor as
the basis for administrative decisions. In particular, terms such as ``too
sweet'', ``insufficiently aromatic'' or ``insufficiently typical'' do not in
themselves constitute evidence of a breach of quality requirements. Declaring
a product non-compliant on the basis of such assessments, without
demonstrating a link to a specific, legally defined criterion or to a
recognized or validated test method, may therefore raise serious doubts as to
the objectivity and legitimacy of the decision
\cite[Art.~34]{Rozporzadzenie2017625}.
